\documentclass[english, parskip=full]{scrartcl}
\usepackage[english]{babel}  % Sprache auf Deutsch 
\usepackage[utf8]{inputenc}  % Kodierung der Datei
\usepackage[left=2cm, right=2cm, top=2cm, bottom=3cm, footskip=1.5cm]{geometry}
\input{myscrarticle}

\newcommand{\TODO}[1]{\color{red} #1 \color{black}}

\newcommand*{\titel}{Various Integral Problems}
\newcommand*{\autor}{Joachim Schwardt}

\hypersetup{pdfauthor={\autor}, pdftitle={\titel}}

%\titlehead{titlehead}
\subject{Mathematical Methods}
\title{\titel}
\author{\autor}
\date{\today}

\begin{document}

%\begin{titlepage}
\maketitle  % Titel setzen
%\tableofcontents  % Inhaltsverzeichnis setzen
%\end{titlepage}

\section{One dimensional}
\subsection{Contour integrals}
\subsubsection{Arc integral}
Consider 
\begin{align}
I_\Gamma &= \intx{\Gamma}{}{ \frac{\powe{\i a x}}{P(x)}}{x}. \label{eqn:arc integral e i a x over P of x}
\end{align}
This integral vanishes for the contour $\Gamma = \klc*{R\powe{\sgn (a) \i t} \,\middle\vert\, t \in \klb*{0,\pi} } $ for any non-constant polynomial $P(x)$ (order $N>0$).
This can be shown by considering the upper bound of the absolute value
\begin{align}
|I_\Gamma| &\le \intx{0}{\pi}{ \frac{1}{|P(R\powe{\i t})|} \abs*{\powe{\i a R\powe{\sgn(a) \i t}}} R }{t} \le \frac{R}{\min |P(R\powe{\i t})|} \intx{0}{\pi}{\powe{- |a| R \sin t}}{t}.
\end{align}
Note that the integral is symmetric about $t=\frac{\pi}{2}$ and that $\sin t \ge \frac{2t}{\pi}$ on the interval $t\in \klb*{0,\frac{\pi}{2}}$.
Also, for sufficiently large $R$ the absolute value of $P$ is governed by the leading order term $P \sim x^N$ where $N > 0$.
Thus we find that there exists some $C>0$ such that
\begin{align}
|I_\Gamma| &\le C \frac{R}{R^N} \intx{0}{\frac{\pi}{2}}{ \powe{-|a| R \frac{2t}{\pi} }}{t} = \frac{C\pi}{2|a|} \frac{1 - \powe{-|a|R}}{R^N} \koml{R\rightarrow\infty}{\longrightarrow} 0.
\end{align}
\subsection{Rational integrals}
\subsubsection{Lorentz integrals}
Consider
\begin{align}
I &= \intr{\frac{1}{1 + x^2}}{x}  = \pi. \label{eqn:1 over 1 plus x squared}
\end{align}
This is solved by substitution $x = \tan t$ since
\begin{align}
I &= 2\intx{0}{\frac{\pi}{2}}{ \frac{1}{1 + \tan^2t} \sec^2t}{t} = 2 \intx{0}{\frac{\pi}{2}}{1 }{t} = \pi.
\end{align}
\subsubsection{Logarithmic Lorentz integral}
For $x_0\in\mathbb{C}$ with $s := \sgn \kla*{\Im x_0} \neq 0$ we have
\begin{align}
I &= \intr{\frac{\log \kla*{x - x_0} }{1 +x^2}}{x} = \frac{\pi}{2} \log \kla*{1 + r^2 + 2|\Im x_0|} - \i \pi \sgn (\Re x_0) \arctan \kla*{\frac{1 + |\Im x_0|}{|\Re x_0|}} . \label{eqn:log x minus x0 over 1 plus x squared}
\end{align}
To avoid residues of logarithms we close the contour in the plane opposite to the location of $x_0$.
The integral over the arc then vanishes (TODO: simple estimate using max value of log for complex inputs, denominator suppresses integral).
The residue theorem then gives
\begin{align}
I &= 2\pi\i \lim_{x \rightarrow -s\i} (x + s\i) \frac{\log (x - x_0)}{1 + x^2} = -s\pi \log \kla*{-s\i - x_0}.
\end{align}
Writing $x_0 = r \powe{\i \phi}$ gives $-s\i - x_0 = R\powe{\i \Phi}$ with $R^2 = (r\cos \phi)^2 + (r|\sin\phi| + 1)^2$ and $\tan\Phi = \frac{-s - r\sin \phi}{-r\cos\phi} = \frac{1 + r|\sin \phi|}{r|\cos\phi|} s \sgn (\cos \phi) $ where we used that $s = \sgn (\sin \phi)$.
The integral then becomes
\begin{align}
I &= -s\pi  \kla*{\log R + \i \Phi} = -s\frac{\pi}{2} \log \kla*{r^2 + 1 + 2r|\sin\phi|} -  \i\pi \arctan \kla*{\frac{1 + r|\sin\phi|}{r|\cos\phi|}} \sgn (\cos\phi) .
\end{align}
\subsubsection{High order logarithmic Lorentz integral}
Consider
\begin{align}
I &= \intr{\frac{\log (1 + x^{4N})}{1 + x^2}}{x} = \TODO{2\pi \sum_{j=1}^N \log \kla*{2 + 2 \abs*{\sin \frac{\pi}{4N}j}} } . \label{eqn:log 1 plus x to N over 1 plus x squared}
\end{align}
\TODO{Result only correct for $N=1$!}\\
Writing $1+x^{4N} = \prod_{j=1}^{4N} \kla*{x - \powe{\i \frac{\pi}{4N}j }} $ gives
\begin{align}
I &= \sum_{j=1}^{4N} I\kla*{x_0 = \powe{\i \frac{\pi}{4N}j}} \kom{(\ref{eqn:log x minus x0 over 1 plus x squared})} \\
&= \frac{\pi}{2} \sum_{j=1}^{4N} \log \kla*{2 + 2 \left\vert \sin \frac{\pi}{4N}j \right\vert} - \i\pi \sum_{j=1}^{4N} \sgn \left(\cos \frac{\pi}{4N}j \right) \arctan\kla*{\frac{1 + |\sin \frac{\pi}{4N}j|}{|\cos \frac{\pi}{4N}j|}} = \\
&= 2\pi \sum_{j=1}^N \log \kla*{2 + 2 \abs*{\sin \frac{\pi}{4N}j}}.
\end{align}
\subsection{Fourier integrals}
\subsubsection{Lorentz integral}
Consider
\begin{align}
I(k) &= \intr{\frac{\powe{\i kx}}{1 + x^2}}{x} = \sgn(k) \pi \powe{-|k|}. \label{eqn:FFT 1 over 1 plus x squared}
\end{align}
This is solved by contour integration.
The poles are $x=\pm\i$ and the integral over the arc described by $x(t) = R\e^{\sgn (k) \i t}$ for $t\in\klb*{0,\pi}$ vanishes in the limit $R\rightarrow\infty$ (follows from \autoref{eqn:arc integral e i a x over P of x}).
The residue theorem then gives
\begin{align}
I(k) &= 2\pi\i \lim_{x\rightarrow \sgn(k)  \i} (x- \sgn(k) \i) \frac{\powe{\i kx}}{1 + x^2} = \frac{2\pi\i}{2\sgn (k) \i} \powe{\i k \sgn (k) \i} = \sgn(k) \pi \powe{-|k|}.
\end{align}


\subsection{Trigonometric integrals}
\subsubsection{Sin integrals}
Consider
\begin{align}
I &= \intr{\frac{\sin\kla*{x - \frac{1}{x}} }{x + \frac{1}{x}}}{x} &= \frac{\pi}{\e^2}. \label{eqn:sin x minus 1 over x over x plus 1 over x}
\end{align}
This is solved by writing $\sin(x) \equiv \Im \powe{\i x}$ and using the residual theorem to get
\begin{align}
I &= 2\pi\i \lim_{z\rightarrow\i} (z - \i) \frac{\powe{\i \kla*{z - \frac{1}{z}}}}{z + \frac{1}{z}} = \pi\i \powe{-2}.
\end{align}
The integral over the arc vanishes due to \autoref{eqn:arc integral e i a x over P of x}.
\subsubsection{Tangent integrals}
Consider
\begin{align}
I &= \intx{0}{\frac{\pi}{2}}{\frac{x}{\tan x}}{x} = \frac{\pi}{2} \log 2. \label{eqn:x over tan x}
\end{align}
This is solved by defining
\begin{align}
I(t) &= \intx{0}{\frac{\pi}{2}}{ \frac{\arctan\kla*{t \tan x} }{\tan x}}{x}.
\end{align}
The derivative is
\begin{align}
I'(t) &= \intx{0}{\frac{\pi}{2}}{ \frac{1}{1 + t^2\tan^2 x} }{x} \kom{u = \tan x} \intx{0}{\infty}{ \frac{1}{1 + t^2u^2} \frac{1}{1+u^2} }{u} = \\
&= \intx{0}{\infty }{\kla*{\frac{\frac{t^2}{t^2-1}}{1 + t^2u^2} - \frac{\frac{1}{t^2-1}}{1 + u^2}} }{u} \kom{(\ref{eqn:1 over 1 plus x squared})} \frac{1}{t^2 - 1} \kla*{\frac{\pi}{2}t - \frac{\pi}{2}} = \frac{\pi}{2} \frac{1}{t + 1}.
\end{align}
Using $I(t=0) = 0$ gives
\begin{align}
I(t) &= I(0) + \intx{0}{t}{I'(s) }{s} =  \frac{\pi}{2} \intx{0}{t}{\frac{1}{s+1}}{s} = \frac{\pi}{2} \log\kla*{t+1}. \label{eqn:arctan t tan x over tan x}
\end{align}
\subsubsection{Bessel integrals}
%https://math.stackexchange.com/questions/4944549/finding-a-general-expression-for-the-improper-integral-int-0-infty-k-1-k2
We have 
%r=1.351; z=0.6699; alpha=0.3153; L=50
%f=lambda x: np.cos(x*z)*special.k0(r*np.sqrt(x**2+alpha**2))
%print(quad(f,0,L), np.pi/2/np.sqrt(r**2+z**2) * np.exp(-alpha*np.sqrt(r**2+z**2)))
\begin{align}
I &= \intx{0}{\infty}{K_0\kla*{r\sqrt{\alpha^2 + x^2}} \cos(xz)}{x} = \frac{\pi}{2} \frac{1}{\sqrt{r^2+z^2}} \powe{-\alpha \sqrt{r^2+z^2}}.
\end{align}
%$K_0(x) = \intx{0}{\infty}{\frac{\cos(xt)}{\sqrt{t^2+1}}}{t}$
To show this, use the integral representation $K_\nu(x) = \frac{1}{2} (x/2)^\nu \intx{0}{\infty}{t^{-\nu-1} \powe{-t-x^2/(4t)}}{t}$ (DLMF 10.32.11, Basset's Integral) to get
\begin{align}
I &= \frac{1}{2} \intx{0}{\infty}{ \frac{\powe{-t}}{t} \intx{0}{\infty}{ \powe{-\frac{r^2}{4t} \kla*{\alpha^2 + x^2}} \cos(xz) }{x}}{t} = \frac{1}{4} \intx{0}{\infty}{ \frac{\powe{-t}}{t} \powe{-\frac{r^2}{4t}\alpha^2} \sqrt{\frac{4t\pi}{r^2}} \powe{-\frac{t}{r^2}z^2} }{t}\\
&= \frac{\sqrt{\pi}}{r} \intx{0}{\infty}{\powe{-u^2 \kla*{1 + \frac{z^2}{r^2}} - \frac{1}{4u^2} r^2\alpha^2}}{u}.
\end{align}
This integral is known ($x= (b/a)^{1/4}y$ and $y= 1/x$ and $u= x-\frac{1}{x} $) %https://math.stackexchange.com/questions/496088/how-to-evaluate-int-0-infty-exp-ax2-frac-bx2-dx-for-a-b0
\begin{align}
J &= \intx{0}{\infty}{\powe{-ax^2 - \frac{b}{x^2}}}{x} = \kla*{\frac{b}{a}}^{1/4} \intx{0}{\infty}{\powe{-\sqrt{ab} \kla*{y^2 + \frac{1}{y^2}}}}{y}\\
&= \kla*{\frac{b}{a}}^{1/4} \intx{1}{\infty}{\powe{-\sqrt{ab} \kla*{y^2 + \frac{1}{y^2}}}}{y} + \kla*{\frac{b}{a}}^{1/4} \intx{1}{\infty}{\powe{-\sqrt{ab} \kla*{\frac{1}{x^2} + x^2}} \frac{1}{x^2}}{x}\\
&= \kla*{\frac{b}{a}}^{1/4} \intx{1}{\infty}{\powe{-\sqrt{ab} \kla*{x^2 + \frac{1}{x^2}}} \kla*{1 + \frac{1}{x^2}} }{x} = \kla*{\frac{b}{a}}^{1/4} \intx{0}{\infty}{\powe{-\sqrt{ab} \kla*{2+u^2}} }{u} = \frac{1}{2} \sqrt{\frac{\pi}{a}} \powe{-2\sqrt{ab}}.
\end{align}
Inserting this gives
\begin{align}
I &= \frac{\pi}{2r} \frac{1}{\sqrt{1 + z^2/r^2}} \powe{-\sqrt{1 + z^2/r^2} \alpha r} = \frac{\pi}{2} \frac{1}{\sqrt{r^2+z^2}} \powe{-\alpha \sqrt{r^2+z^2}}.
\end{align}
%Or use $K_\nu(xy) = \frac{1}{\sqrt{\pi}} \Gamma\kla*{\nu + \frac{1}{2}} (2y/x)^\nu \intx{0}{\infty}{\kla*{t^2+y^2}^{-\nu-1/2} \cos(xt)}{t}$ to get
%\begin{align}
%I &= \intx{0}{\infty}{\cos(rt) \intx{0}{\infty}{ \frac{\cos(xz)}{\sqrt{t^2+\alpha^2+x^2}} }{x}}{t} = 
%\end{align}%\intx{0}{\infty}{\cos(rt) K_0\kla*{z \sqrt{t^2+\alpha^2} }}{t}
%\frac{1}{4} \intr{ \intr{ \frac{\powe{\i rt + \i xz}}{\sqrt{\alpha^2+t^2+x^2}}}{x} }{t}\\
%&= \frac{1}{4} \intx{0}{2\pi}{\intx{0}{\infty}{\frac{\powe{\i\kla*{r\cos\phi + z\sin\phi} R}}{\sqrt{\alpha^2+R^2}} R}{R}}{\phi}
%Or note that $\intx{0}{\infty}{t^{\mu-1} K_\nu(t)}{t} = 2^{\mu-2} \Gamma\kla*{\frac{\mu-\nu }{2}} \Gamma\kla*{\frac{\mu+\nu }{2}}$ and expand,
%\begin{align}
%I &=
%\end{align}

Next, let us consider the similar integral
\begin{align}
I &= \intx{0}{\infty}{K_1\kla*{r\sqrt{\alpha^2 + x^2}} \sin(xz)}{x}.
\end{align}
Using the same integral representation, we have%\footnote{$\intx{0}{\infty}{ \powe{-a^2x^2} x^n}{x} = \frac{1}{2a^{n+1}} \Gamma\kla*{\frac{n+1}{2}}$}
%r=1.351; z=0.6699; alpha=0.9153; L=50
%f=lambda x: np.sin(x*z)*special.k1(r*np.sqrt(x**2+alpha**2))
%print(quad(f,0,L))
%quad(lambda t: r/4 * np.exp(-t)/t**2 * quad(lambda x: np.sqrt(alpha**2+x**2) * np.sin(x*z) * np.exp(-r**2/4/t * (alpha**2+x**2)), 0, L)[0], 0, L)
\begin{align}
I &= \frac{r}{4} \intx{0}{\infty}{ \frac{\powe{-t}}{t^2} \intx{0}{\infty}{ \powe{-\frac{r^2}{4t} \kla*{\alpha^2 + x^2}} \sqrt{\alpha^2 + x^2} \sin(xz)}{x}}{t} \\
&= \frac{r}{4} \intx{0}{\infty}{ \frac{\powe{-t}}{t^2} \intx{\alpha}{\infty}{ \powe{-\frac{r^2}{4t} u^2} u  \sin\kla*{z \sqrt{u^2-\alpha^2} } \frac{u}{\sqrt{u^2-\alpha^2}} }{u}}{t}\\
&= \frac{r}{4} \sum_{n\ge 0} \frac{(-1)^n z^{2n+1}}{(2n+1)!} \intx{0}{\infty}{ \frac{\powe{-t}}{t^2} \intx{\alpha}{\infty}{ \powe{-\frac{r^2}{4t} u^2} u^2 \kla*{u^2-\alpha^2}^n }{u}}{t}
\end{align}
Or using $K_1 = -K_0'$, we may write%https://math.stackexchange.com/questions/2204475/derivative-of-bessel-function-of-second-kind-zero-order
\begin{align}
I &= -\partial_r \intx{0}{\infty}{K_0\kla*{r\sqrt{\alpha^2 + x^2}} \frac{\sin(xz)}{\sqrt{\alpha^2 + x^2}}}{x}
\end{align}
Or consider an expansion in $\alpha$.
For example, we have
\begin{align}
I(\alpha=0) &= \intx{0}{\infty}{K_1(rx) \sin(xz)}{x} = \frac{r}{4} \intx{0}{\infty}{ \frac{\powe{-t}}{t^2} \intx{0}{\infty}{ x \powe{-\frac{r^2}{4t}x^2} \sin(xz) }{x} }{t}\\
&= - \partial_z \frac{r}{8} \intx{0}{\infty}{ \frac{\powe{-t}}{t^2} \intr{ \powe{-\frac{r^2}{4t}x^2} \cos(xz) }{x} }{t} = -\partial_z\frac{r}{8} \intx{0}{\infty}{\frac{\powe{-t}}{t^2} \sqrt{\frac{4\pi t}{r^2}} \powe{- \frac{z^2}{r^2} t} }{t}\\
&= -\partial_z\frac{\sqrt{\pi}}{4} \sqrt{1 + \frac{z^2}{r^2}} \Gamma\kla*{-\frac{1}{2}} = \frac{\pi}{2r} \frac{z}{\sqrt{r^2+z^2}}.
\end{align}
Similarly, note that
\begin{align}
\partial_{\alpha^2}I(\alpha=0) &= \intx{0}{\infty}{ K_1'(rx) \frac{1}{2x} \sin(xz) }{x}???
\end{align}

\subsection{Exponential integrals}
\subsubsection{Gamma function}
The Gamma function is
\begin{align}
\Gamma(x) &= \intx{0}{\infty}{t^{x-1} \powe{-t}}{t}. \label{eqn:gamma function}
\end{align}
\subsubsection{Beta function}
The Beta function is
\begin{align}
B(x,y) &= \intx{0}{1}{t^{x-1} (1 - t)^{y-1}}{t} = 2\intx{0}{\frac{\pi}{2}}{\sin^{2x-1}t \cos^{2y-1} t}{t} =  \frac{\Gamma(x)\Gamma(y)}{\Gamma(x+y)}. \label{eqn:beta function}
\end{align}
The first equality is found by $t = \sin^2 \phi$.
To derive the relation to the Gamma function consider
\begin{align}
\Gamma(x) \Gamma(y) &= \intx{0}{\infty}{ \intx{0}{\infty}{ t^{x-1} s^{y-1} \powe{-(s+t)}}{s}}{t} \kom{s = u - t} \\
&= \intx{0}{\infty}{\intx{t}{\infty}{ \powe{-u} t^{x-1} (u - t)^{y-1} }{u}}{t} \kom{u = ts} \\
&= \intx{0}{\infty}{ \intx{1}{\infty}{ \powe{-ts} t^{x-1} t^{y-1} (s - 1)^{y-1} t}{s} }{t} \kom{u = ts} \\
&= \intx{0}{\infty}{ \intx{1}{\infty}{ \powe{-u} u^{x+y-1} s^{1-x-y} (s-1)^{y-1} \frac{1}{s} }{s} }{u} \kom{t = \frac{1}{s}} \\
&= \Gamma(x+y) \intx{1}{0}{ t^{x+y-1} \kla*{\frac{1}{t} - 1}^{y-1} t \frac{-1}{t^2} }{t} = \Gamma(x+y) B(x,y).
\end{align}
\subsubsection{Gauss integral}
Consider
\begin{align}
I &= \intr{\powe{-x^2}}{x} = \sqrt{\pi}. \label{eqn:e to the minus x squared}
\end{align}
Note that
\begin{align}
I^2 &= \intr{\intr{ \powe{-(x^2+y^2)}}{x}}{y} = \intx{0}{1}{\intx{0}{2\pi}{ \powe{-r^2} r}{\phi}}{r} = 2\pi \frac{-1}{2} \powe{-r^2}\Big\vert_{r=0}^{r=\infty} = \pi.
\end{align}
\subsubsection{General gauss integral}
Consider
\begin{align}
I &= \intr{\powe{-ax^2 + bx}}{x} = \sqrt{\frac{\pi}{a}} \powe{\frac{b^2}{4a}}. \label{eqn:e to the minus a x squared plus b x}
\end{align}
This is derived by the substitution $x = y + \frac{b}{2a}$ and using \autoref{eqn:e to the minus x squared}
\begin{align}
I &= \intr{\powe{-ay^2 + \frac{b^2}{4a}}}{y}.
\end{align}
\subsubsection{Stirling's formula}
Consider the integral
\begin{align}
\frac{\Gamma(x)}{x^{x}} &= \intx{0}{\infty}{t^{x-1}\powe{-tx}}{t} = \intx{0}{\infty}{ \powe{(x-1)\log(t) - tx} }{t} = \intr{\powe{xu - x\powe{u}}}{u}.
\end{align}
Using $f(u)=\powe{u}-u = 1 + \frac{u^2}{2} + g(u)$ with $g(u) = \sum_{m\ge 3} \frac{u^m}{m!}$ yields
\begin{align}
\frac{\Gamma(x)}{x^{x}} &= \powe{-x} \intr{\powe{-x\frac{u^2}{2} - xg(u)}}{u} = \powe{-x} \sum_{n\ge 0} \frac{x^n}{n!} \intr{g(u)^n \powe{-\frac{x}{2}u^2}}{u} \\
&= \powe{-x} \frac{1}{\sqrt{x}} \sum_{n\ge 0} \frac{x^n}{n!} \intr{g\kla*{\frac{t}{\sqrt{x}}}^n \powe{-\frac{1}{2}t^2}}{t}.
\end{align}
Rearranging gives the series expansion
\begin{align}
\powe{x} x^{-x} \frac{1}{\sqrt{2\pi x}} x! &= \frac{1}{\sqrt{2\pi}} \sum_{n\ge 0} \sum_{m_1\ge 3} \dots \sum_{m_n\ge 3} \frac{x^{-(m_1+\dots+m_n)/2}}{m_1!\dots m_n!} \intr{t^{m_1+\dots+m_n} \powe{-\frac{1}{2}t^2}}{t}\\
&= \frac{1}{\sqrt{\pi}} \sum_{n\ge 0} \sum_{\textbf{m}\ge 3, |\textbf{m}| \text{ even}} \frac{1}{\textbf{m}!} \Gamma\kla*{\frac{|\textbf{m}|+1}{2}} \frac{1}{(2x)^{|\textbf{m}|/2}}.
\end{align}
The leading order of the RHS expansion is simply 1, and since $|\textbf{m}|=2$ is impossible, the next term is $\frac{3}{4} \frac{1}{(2x)^2} = \frac{3}{16x^2}$.
TODO something is wrong here
\subsubsection{Exponential cosine gauss integral}
\url{https://math.stackexchange.com/questions/4933804/exponential-cosine-gaussian-integral}
\begin{align}
I &= \intr{\powe{-ax^2 + b\cos(x)}}{x}.
\end{align}
\subsubsection{Super-exponential integrals}
Consider
\begin{align}
I &= \intx{0}{1}{x^{px}}{x} = \frac{-1}{p}\sum_{n\ge 1} \frac{(-p)^n}{n^n}. \label{eqn:x to the p x}
\end{align}
This can be solved by rewriting $x^{px} = \powe{px\log x}$ and substituting $u = -\log x$.
Expanding the outer exponential into a Taylor series then yields
\begin{align}
I &= \intx{0}{\infty}{\powe{-pu \powe{-u}} \powe{-u} }{u} = \sum_{n\ge 0} \frac{(-p)^n}{n!} \intx{0}{\infty}{u^{n} \powe{-(n+1)u}}{u} = \sum_{n\ge 0} \frac{(-p)^n}{n!} \frac{\Gamma (n+1)}{(n+1)^{n+1}}.
\end{align}
\subsection{Exotic integrals}
\subsubsection{Exhaustive method for functions of bounded variation}
If $f$ has bounded variation then
\begin{align}
I &= \intx{a}{b}{ f(x)}{x} = \sum_{n\ge 1} \frac{1}{2^n} \sum_{k=1}^{2^n-1} (-1)^{k+1} f \kla*{a + \frac{k}{2^n}(b-a)}. \label{eqn:exhaustive method f of x from a to b}
\end{align}
For a simple example let $a=0$, $b=1$ and $f(x) = x^2$.
Then\footnote{Note that $\sum_{k=1}^n (-1)^{k+1} k^2 = (-1)^{n+1} \frac{n(n+1)}{2} $.}
\begin{align}
I &= \sum_{n\ge 1} \frac{1}{2^n} \sum_{k=1}^{2^n - 1} (-1)^{k+1} \frac{k^2}{2^{2n}} = \sum_{n\ge 1} \frac{1}{2^{3n}} (-1)^{2^n} \frac{(2^n - 1) 2^n}{2} = \\
&= \frac{1}{2} \sum_{n\ge 1} \frac{1}{2^n} - \frac{1}{2} \sum_{n\ge 1} \frac{1}{2^{2n}} = \frac{1}{2} - \frac{1}{2} \frac{1}{3} = \frac{1}{3}.
\end{align}
For a more complicated example let $a=0$, $b=\pi$ and $f(x) = \sin(x)$.
Then\footnote{Note that $\sum_{k=1}^n (-1)^{k+1} \sin \kla*{\frac{k\pi}{n-1}} = \Im \frac{\powe{\i \pi / n} + (-1)^{n+1}}{\powe{\i \pi / n} + 1} = \begin{cases} \tan\kla*{\frac{\pi}{2n}} & n\ \Text{even} \\ 0 & n\ \Text{odd} \end{cases}$.}
\begin{align}
I &= \pi \sum_{n\ge 1} \frac{1}{2^n} \sum_{k=1}^{2^n - 1} (-1)^{k+1} \sin\kla*{\frac{k\pi}{2^n} } = \pi\sum_{n\ge 1} \frac{1}{2^n} \tan\kla*{\frac{\pi}{2^{n+1}}} \TODO{=2} .
\end{align}
%Also note that the result can be rewritten.
%Let the summand expression be $S(n,k)$.
%Then
%\begin{align}
%I &= \sum_{n\ge 1} \sum_{k=1}^{2^n - 1} S(n,k) = \sum_{k\ge 1} \sum_{n\ge \lfloor \log_2(k+1) \rfloor} S(n,k) = \sum_{k\ge 1} \sum_{n\ge 0} S\kla*{n + \lfloor \log_2(k+1) \rfloor, k} = \\
%&= \sum_{k\ge 1} (-1)^{k+1} \sum_{n\ge 0} \frac{1}{2^n 2^{\lfloor \log_2(k+1) \rfloor} } f \kla*{a + \frac{k}{2^n 2^{\lfloor \log_2(k+1) \rfloor} } (b-a)}.
%\end{align}
\subsubsection{Sin integrals}
Consider
\begin{align}
I &= \intx{0}{1}{\frac{\sin(\pi x)}{x^x (1-x)^{1-x}}}{x} = \frac{\pi}{\e}. \label{eqn:sin pi x over x to x times 1 minus x to 1 minus x}
\end{align}
\TODO{How?}
%%Rewriting gives
%%\begin{align}
%%I &= \intx{0}{1}{\frac{\sin (\pi x)}{1-x} \kla*{\frac{1}{x} - 1}^x }{x}
%%\end{align}
%%what is 
%%\begin{align}
%%x\log x + (1-x)\log (1-x) &= 
%%\end{align}
\subsubsection{Hardy-Ramanujan integral partition function}
Show that %https://math.stackexchange.com/questions/5058534/show-that-int-0-infty-frac1ueu-1-frac1u2-frace-u2u-du
\begin{align}
I &= \intx{0}{\infty}{\kla*{\frac{1}{x\kla*{\powe{x}-1}} - \frac{1}{x^2} + \frac{\powe{x}}{2x}}}{x} = -\frac{1}{2} \log \kla*{2\pi}.
\end{align}
%Using the series expansion %https://en.wikipedia.org/wiki/Bernoulli_number#Generating_function
%\begin{align}
%\frac{x}{\powe{x}-1} &= \sum_{n\ge 0} \frac{B_n^-}{n!} x^n,
%\end{align}
%where $B_n^- = $
Note the series expansion %https://proofwiki.org/wiki/Power_Series_Expansion_for_Hyperbolic_Cotangent_Function
\begin{align}
\coth(x) &= \sum_{n\ge 0} \frac{2^{2n}}{(2n)!} B_{2n} x^{2n-1}
\end{align}
where $B_{2n} = (-1)^{n+1} \frac{2(2n)!}{(2\pi)^{2n}} \zeta(2n)$ are the even Bernoulli numbers and $\zeta(s)$ is the Riemann Zeta function.
We also have the related series ($|x|<2\pi$) %https://proofwiki.org/wiki/Definition:Bernoulli_Numbers/Archaic_Form
\begin{align}
\frac{x}{\powe{x}-1} &= 1 - \frac{x}{2} + \sum_{n\ge 1} \frac{B_{2n}}{(2n)!} x^{2n}
\end{align}
Using this, the finite part of the integral can be written as
\begin{align}
I_1 &= \intx{0}{2\pi}{\kla[\bigg]{\frac{\powe{x} - 1}{2x} + \sum_{n\ge 0} \frac{B_{2n+2}}{(2n+2)!} x^{2n}}}{x}\\
&= \sum_{n\ge 0} \intx{0}{2\pi}{\kla*{\frac{1}{2} \frac{x^n}{(n+1)!} + \frac{B_{2n+2}}{(2n+2)!} x^{2n}}}{x}\\
&= \sum_{n\ge 0} \kla*{\frac{1}{2n} \frac{(2\pi)^{n+1}}{(n+1)!} + \frac{B_{2n+2}}{2n+1} \frac{(2\pi)^{2n+1}}{(2n+2)!} }
\end{align}
\section{Two dimensional}
\subsection{Trigonometric integrals}
\subsubsection{Light-cone integration}
Consider, for $\Re z < b$,
\begin{align}
I(z) &= \intx{0}{\infty}{\powe{xz} \csch(x)^b}{x} = 2^b \intx{0}{\infty}{\powe{xz} \kla*{1 - \powe{-2x}}^{-b} \powe{-bx}}{x}\\
&= 2^{b-1} \intx{0}{1}{u^{\frac{b-z}{2} - 1} \kla*{1 - u}^{-b}}{u} = 2^{b-1} B\kla*{1-b, \frac{b+z}{2}} = 2^{b-1} \Gamma(1-b) \frac{\Gamma\kla*{\frac{b}{2} - \frac{z}{2}}}{\Gamma\kla*{1 - \frac{b}{2} - \frac{z}{2}}}.
\end{align}
\subsubsection{$J$-integration problem}
Consider
\begin{align}
J_{b,0}(k,\omega) &= \intx{0}{\pi}{\intr{\powe{\i kx - \i 2n \tau} \csc(\tau + \i x)^b \csc(\tau - \i x)^b}{x} }{\tau} \\
&= 2^{2b-2} \Gamma(1-b)^2 \sin(\pi b) \frac{\Gamma\kla*{\frac{b}{2} - \i \frac{k}{4} + \frac{n}{2}}}{\Gamma\kla*{1 - \frac{b}{2} - \i \frac{k}{4} + \frac{n}{2}}} \frac{\Gamma\kla*{\frac{b}{2} + \i \frac{k}{4} + \frac{n}{2}}}{\Gamma\kla*{1 - \frac{b}{2} + \i \frac{k}{4} + \frac{n}{2}}}
\end{align}
%b=0.431; k=0.4331; omega=2j; L=3
%_j_bn(k, omega, b), c_quad(lambda tau: np.exp(-omega*tau) * c_quad(lambda x: np.exp(1j*k*x)/np.sin(tau+1j*x)**b/np.sin(tau-1j*x)**b, -L, L)[0], 0, np.pi)
%prefactor = np.pi * 2**(2*b-2) * _gamma_real(1-b) / _gamma_real(b)
%gamma_p = _gamma_rel(b/2 - 1j*(omega+k)/4, 1 - b/2 - 1j*(omega+k)/4)
%gamma_m = _gamma_rel(b/2 - 1j*(omega-k)/4, 1 - b/2 - 1j*(omega-k)/4)
%prefactor * gamma_p * gamma_m
where $\omega=\i 2n$ for $n\in\mathbb{Z}$.
Note that $k_-z_+ + k_+z_- = \frac{\i k- 2n}{2} \kla*{x +\i \tau} + \frac{\i k+ 2n}{2} \kla*{x -\i \tau} = \i kx - \i 2n\tau$ where $z_\pm = x \pm \i \tau$ and $k_\pm = \frac{\i k \pm 2n}{2}$.
Using $\csc(\i x) = \frac{1}{2\i} \kla*{\powe{-x} - \powe{x}} = \i \csch(x)$ we then have $\csc(\tau \pm \i x) = \csc(\pm \i z_\mp) = \pm \i \csch(z_\mp)$.
Naive substitution would then give
\begin{align}
J_{b,0}^\text{naive sub.} &= \intx{?}{?}{ \intx{?}{?}{ \powe{k_-z_+} \powe{k_+z_-} \csch(z_+)^b \csch(z_-)^b }{z_+}}{z_-} \overset{?}{=} I(-k_+) I(k_-) = \csc(\pi b) J_{b,0}.
\end{align}
This is surprisingly close to the correct answer, but of course we are missing the $\sin(\pi b)$-factor and it is unclear how to make this multidimensional complex substitution rigourous.
In particular, the integration region has to somehow become $[0,\infty]$ and $[-\infty, 0]$ for $z_+$ and $z_-$ respectively.

Is there an alternative way to derive $J_{b,0}$ then?
Let us try a brute-force approach (for $n\rightarrow-n$)
\begin{align}
J_{b,0} &= 2^b\intx{0}{\pi}{\intr{\powe{\i kx - \i 2n\tau} \klb*{\cosh(2x) - \cos(2\tau)}^{-b}}{x} }{\tau} \\
&= 4^b \int_{|z|=1}\intx{0}{\infty}{u^\frac{\i k}{2} z^n \klb*{u + \frac{1}{u} - z - \frac{1}{z}}^{-b} \frac{1}{2u} }{u} \frac{1}{2\i z} \,\d z\\
&= 2^{2b-2} \frac{1}{\i} \intx{0}{\infty}{ u^{\frac{\i k}{2} + b - 1} \int_{|z|=1} z^{n+b-1} \klb*{(u^2+1)z - z^2u - u}^{-b}\,\d z }{u}\\
&= 
\end{align}
\section{Multi dimensional}
\subsection{Exponential integrals}
\subsubsection{Gaussian integral}
Consider
\begin{align}
I &= \intrnd{N}{\powe{-\brc*{x}{A}{x} + \brb*{b}{x} }}{x} = \sqrt{\frac{\pi^N}{\det A}} \powe{\frac{1}{4} \brc*{b}{A^{-1}}{b} }. \label{eqn:e to the minus x A x plus b x}
\end{align}
This result holds for any symmetric matrix $A=UDU^T$. 
\TODO{maybe we also have to assume positive definiteness for positive eigenvalues?}\\
To prove it substitute $y = U^Tx$ to get
\begin{align}
I &= \intrnd{N}{ \powe{-\brc*{y}{D}{y} + \brc*{b}{U}{y} }}{y} = \prod_{j=1}^N \intr{\powe{-\lambda_j y^2 + (U^Tb)_j y}}{y} = \\
&= \prod_{j=1}^N \sqrt{\frac{\pi}{\lambda_j}}  \powe{\frac{(U^Tb)_j^2}{4\lambda_j}} = \sqrt{\frac{\pi^N}{\det A}} \powe{\frac{1}{4} \sum_{j=1}^N (U^T b)_j \frac{1}{\lambda_j} (U^T b)_j } = \\
&= \sqrt{\frac{\pi^N}{\det A}} \powe{\frac{1}{4} \brc*{b}{UD^{-1} U^T}{b} } = \sqrt{\frac{\pi^N}{\det A}} \powe{\frac{1}{4} \brc*{b}{A^{-1}}{b} }.
\end{align}
\subsubsection{Gaussian complex integral}
Consider
\begin{align}
I &= \int_{\mathbb{C}^N}^{} \powe{-\brc*{x}{A}{x} } \,\mathrm{d}^Nx\mathrm{d}^Nx^\ast = \frac{\pi^N}{\det A}. \label{eqn:e to the minus x ast A x}
\end{align}
Assuming a hermitian matrix with $A=UDU^\dagger$ we can write $\brc*{x}{A}{x } = \brc*{y}{D}{y}$ with $x = Uy$.
Then the integral becomes
\begin{align}
I &= \prod_{j=1}^N \intr{\intr{ \powe{-\lambda_j y^\ast y} }{y^\ast}}{y} = \prod_{j=1}^N \intr{\intr{ \powe{-\lambda_j (y_r^2 + y_i^2)}}{y_r}}{y_i} = \frac{\pi^N}{\det A}.
\end{align}
\subsubsection{Gaussian complex integral matrix elements}
Consider
\begin{align}
I &= \frac{\int_{\mathbb{C}^N}^{} \powe{-\brc*{x}{A}{x} } x_k^\ast x_l \,\mathrm{d}^Nx\mathrm{d}^Nx^\ast}{\int_{\mathbb{C}^N}^{} \powe{-\brc*{x}{A}{x} } \,\mathrm{d}^Nx\mathrm{d}^Nx^\ast} = (A^{-1})_{lk}. \label{eqn:e to the minus x ast A x xj xk normalized}
\end{align}
Differentiating \autoref{eqn:e to the minus x ast A x} with respect to $A_{kl}$ gives a factor of $-x_k^\ast x_l$ in the integrand.
Thus the integral is given by
\begin{align}
I &= -\frac{\det A}{\pi^N} \frac{\partial }{\partial A_{kl}} \frac{\pi^N}{ \det A} = \det A \frac{1}{(\det A)^2} \frac{\partial }{\partial A_{kl}} \det A = (A^{-1})_{lk}.
\end{align}
The last step involves the observation
\begin{align}
\frac{\partial}{\partial A_{jk}} \log \det A &= \frac{\partial}{\partial A_{jk}} \Tr \log A = \Tr A^{-1} \frac{\partial}{\partial A_{jk}} A = \\
&= \Tr \kla*{(A^{-1})_{mn} (\delta_{jn} \delta_{kl})_{nl} }_{ml} = \sum_{mn} (A^{-1})_{mn} \delta_{jn} \delta_{km} = (A^{-1})_{kj}.
\end{align}
\section{Convolutions}
We define the convolution of two functions as
\begin{align}
\kla*{f\ast g}(x) &= \intx{ }{ }{f(y)g(x-y)}{y}.
\end{align}
This operation is symmetric, and under Fourier transform reduces to simple multiplication.
Explicitly, if $f(k) = \intx{ }{ }{\powe{-\i kx} f(x)}{x}$, then
\begin{align}
\kla*{f\ast g}(k) &= \intx{ }{ }{\powe{-\i kx} \intx{ }{ }{f(y) g(x-y)}{y} }{x} = \intx{ }{ }{\intx{ }{ }{\powe{-\i k(x+y)} f(y) g(x)}{y}}{x} = f(k)g(k).
\end{align}
This works so long as the integration domain is either infinite, or the functions satisfy periodic translational invariance, such that the substitution does not make the boundaries depend on the other integration variable.
\subsection{Stable functions}
A function is stable under convolution if $f_a \ast f_b = f_c$ for some $c=c(a,b)$, i.e. if the convolution of a function from a certain parameter family with another results in itself.
A special case is $f_a(k)=\klb*{f(k)}^a$, because (assuming convergence) we then find
\begin{align}
\kla*{f_a\ast f_b}(x) &= \intx{ }{ }{ \frac{1}{\mathcal{N}}\intx{ }{ }{\powe{\i ky} \klb*{f(k)}^a}{k} \frac{1}{\mathcal{N}}\intx{ }{ }{\powe{\i k'(x-y)} \klb*{f(k')}^b}{k'} }{y}\\
&= \frac{1}{\mathcal{N}} \intx{ }{ }{\powe{\i kx} \klb*{f(k)}^{a+b}}{k} = f_{a+b}(x) ,
\end{align}
where $\mathcal{N}$ is the normalization factor in the Fourier Transform (e.g. $2\pi$ for $\mathbb{R}$) and $\frac{1}{\mathcal{N}} \intx{ }{ }{\powe{\i kx}}{k} = \delta(x)$ holds (assuming such an identity is possible for the integration domain, but otherwise FTs are ill-defined).
Similarly, $f_a(k) = \intx{ }{ }{\powe{-\i kx} \klb*{f(x)}^a}{x}$ also leads to
\begin{align}
\kla*{f_a \ast f_b}(k) &= \intx{ }{ }{ \intx{ }{ }{\powe{-\i k'x} \klb*{f(x)}^a}{x} \intx{ }{ }{\powe{-\i (k-k')y} \klb*{f(y)}^b}{y} }{k'}\\
&= \mathcal{N} \intx{ }{ }{\powe{-\i kx} \klb*{f(x)}^{a+b}}{x} = \mathcal{N} f_{a+b}(k).
\end{align}
To summarize this result, any function (with fineprint regarding convergence and normalizations) of the form $f(x) = \frac{1}{\mathcal{N}} \intx{ }{ }{\powe{\i kx} \klb*{f(k)}^a}{k}$ with some generator $f(k)$ is stable.
\subsubsection{Cosecant generator}
The generator $f(x) = \csc\kla*{\tau + \i x}$ on $x\in\mathbb{R}$ leads to\footnote{$\intx{0}{\infty}{\powe{\i 2kx} \csch(x)^a}{x} = 2^{a-1} B\kla*{1-a, \frac{a}{2} - \i k}$}
%from scipy import special
%from scipy.integrate import quad
%g=special.gamma
%beta=lambda x,y: g(x)*g(y)/g(x+y)
%cquad = lambda *args: quad(*args, complex_func=True)
%tau=0.735; a=0.351; k=1.555; L=200
%f=lambda x: 1/np.sin(tau + 1j*x)
%Fi=lambda k,a: cquad(lambda x: np.exp(-1j*k*x) * f(x)**a, -L, L)
%F=lambda k,a: 2**(a-1) * beta((a+1j*k)/2, (a-1j*k)/2) * np.exp(k*(tau-np.pi/2))
%Fi(k,a), F(k,a)
\begin{align}%see integral_sandbox Master LL cosecant power 2D Fourier for details
f_a(k) &= \intr{\powe{-\i kx} \csc\kla*{\tau + \i x}^a}{k} = (2\i)^a \powe{-\i\tau a} \intr{\powe{-\i kx} \powe{-ax} \kla*{\powe{-2x} - \powe{-2\i\tau}}^{-a}}{k}\\
&= \frac{2^a}{\Gamma(a)} \powe{-\i\tau a} \powe{\i \frac{\pi}{2}a}\intx{0}{\infty}{v^{a-1} \powe{v\powe{-2\i\tau}} \intr{\powe{-x(a+\i x)} \powe{-v \powe{-2x}}}{x}}{v}\\
&= \dots = 2^{a-1} B\kla*{\frac{a+\i k}{2}, \frac{a-\i k}{2}} \powe{k\kla*{\tau - \frac{\pi}{2}}},
\end{align}
%b=0.661
%cquad(lambda kp: F(kp,a) * F(k-kp,b), -L,L), 2*np.pi* F(k,a+b)
and therefore $f_a\ast f_b = 2\pi f_{a+b}$.
%f_a = lambda x,a: 1/(2*np.pi) * 2**(a-1)/g(a) * g((a+1j*x)/2) * g((a-1j*x)/2) * np.exp(x*(tau-np.pi/2))
%cquad(lambda xp: f_a(xp,a) * f_a(x-xp,b), -L,L), f_a(x,a+b)
Written differently, we have that $f_a(x) = \frac{1}{2\pi} \frac{2^{a-1}}{\Gamma(a)} \Gamma\kla*{\frac{a+\i x}{2}} \Gamma\kla*{\frac{a-\i x}{2}} \powe{x\kla*{\tau - \frac{\pi}{2}}}$ satisfies $f_a\ast f_b = f_{a+b}$.


%\begin{thebibliography}{9}
%\bibitem{example} description, \url{https://www.exmapleurl}, day.\,month~2021.
%\end{thebibliography}

\end{document}